\section{Full four-terminal forest sectors force planarity}
\label{sec:path-six-four-sector-planarity}

The distinction between arbitrary terminal-state vectors and graph states is essential for six-edge cuts.
The next theorem supplies a geometric obstruction to simultaneous saturation of all four-terminal forest sectors.

\begin{theorem}[Planarity from full four-terminal sectors]
\label{thm:path-six-full-four-planar}
Let $X$ be a connected loopless graph of maximum degree three.
Let $p_1,\ldots,p_k$ be distinct degree-two ports, where $k\ge6$.
Suppose every edge set with odd-degree vertex set equal to any four selected ports is a forest.
Then $X$ is planar.

Equivalently, for a connected normalized $k$-pole of rank $h$, the hypothesis says that every four-terminal forest total equals $2^h$.
\end{theorem}

\begin{proof}
Suppose $X$ is nonplanar.
Kuratowski's theorem supplies a subdivision of $K_5$ or $K_{3,3}$ as a subgraph.
The maximum degree excludes a subdivision of $K_5$.
Thus $X$ contains a subdivision $K$ of $K_{3,3}$.
Call its nine subdivided edge paths the core paths.

Each branch vertex of $K$ already has three incident edges in $K$.
It cannot be a port or have a neighbor outside $K$.
Every port lying in $K$ is therefore internal to one core path.
Assign that port to its core path.

Each component of $X-V(K)$ has an edge to $K$, since $X$ is connected.
Choose one such attachment edge for each component.
Its endpoint in $K$ is internal to a core path, by the degree-three condition at branch vertices.
Assign every port in that component to the core path containing this chosen attachment.
This assigns all $k$ ports to the nine core paths.

For each core edge $e$, let $w_e$ count the ports assigned to its path.
The graph $K_{3,3}$ has nine four-cycles.
Each core edge occurs in four of these cycles and is omitted from five.
Consequently, the average total assigned weight on edges omitted from a four-cycle is
\[
 \frac59\sum_e w_e=\frac{5k}{9}\ge\frac{10}{3}>3.
\]
Choose a four-cycle whose omitted edges have total assigned weight at least four.
Let $C$ be its subdivided cycle in $K$.

The two branch vertices outside this four-cycle remain joined by their core path after deleting $V(C)$.
Each of the other four omitted core paths has one of these two branch vertices as an endpoint.
Its surviving interior remains connected to that endpoint.
Thus all surviving interiors of the five omitted core paths lie in one component of $K-V(C)$.

Every port assigned to an omitted core path survives in the same component of $X-V(C)$.
For a port inside $K$, this follows from the preceding paragraph.
For a port outside $K$, its entire outside component survives, and its chosen attachment joins it to that surviving core path.
Hence one component of $X-V(C)$ contains at least four ports.

Choose four such ports and call their set $S$.
A connected graph admits an edge set with odd-degree vertex set exactly $S$, since $|S|$ is even.
For example, use the unique parity solution on a spanning tree of that component.
This edge set is disjoint from $C$.
Adding $E(C)$ gives an edge set in $X$ with odd-degree set exactly $S$ that contains a cycle.
This contradicts the stated forest hypothesis.
Therefore $X$ is planar.
\end{proof}

\begin{corollary}[A missing forest state in every six-pole of girth at least seven]
\label{cor:path-six-girth-seven-missing-state}
Let $X$ be a normalized six-pole of rank $h$ and girth at least seven.
Then at least one four-terminal forest total is at most $2^h-1$.
\end{corollary}

\begin{proof}
Suppose every four-terminal forest total equals $2^h$.
Theorem~\ref{thm:path-six-full-four-planar} makes $X$ planar.
Write $n=|V(X)|$ and $e=|E(X)|$.
Its degrees give $2e=3n-6$.

The graph has minimum degree two and contains a cycle.
In a planar embedding, every face boundary walk has length at least seven.
Indeed, deleting bridge excursions from a face boundary leaves at least one cycle, whose length is at least the girth.
Euler's formula and the sum of face-boundary lengths yield
\[
 2e\ge7(e-n+2),\qquad
 2e\le\frac{14}{5}(n-2).
\]
Combining this with the degree sum gives $n\le2$.
This is impossible because the six ports are distinct.
Thus some forest total is strictly smaller than $2^h$.
Its integrality gives the stated bound.
\end{proof}

The proof requires saturation only for four-terminal forest sectors.
It does not require saturation of the six-terminal sector.
The missing-state gap is exactly at least $2^{-h}$ after normalization.
This gap is useful when the opposite six-cut shore has unbounded rank.
